% !TEX program = xelatex
\documentclass[10pt,letterpaper,twocolumn]{article}

\usepackage[top=0.85in, bottom=0.85in, left=0.75in, right=0.75in]{geometry}
\usepackage{fontspec}
\usepackage{unicode-math}
\setmainfont{Latin Modern Roman}[
  Extension=.otf,
  UprightFont=lmroman10-regular,
  BoldFont=lmroman10-bold,
  ItalicFont=lmroman10-italic,
  BoldItalicFont=lmroman10-bolditalic,
]
\setsansfont{Latin Modern Sans}[
  Extension=.otf,
  UprightFont=lmsans10-regular,
  BoldFont=lmsans10-bold,
  ItalicFont=lmsans10-oblique,
  BoldItalicFont=lmsans10-boldoblique,
]
\newfontfamily\acbold{ywft-processing-bold}[Path=../../system/public/type/webfonts/,Extension=.ttf]
\newfontfamily\aclight{ywft-processing-light}[Path=../../system/public/type/webfonts/,Extension=.ttf]
\setmonofont{Latin Modern Mono}[
  Extension=.otf,
  UprightFont=lmmono10-regular,
  ItalicFont=lmmono10-italic,
  Scale=0.85,
]

\usepackage{xcolor}
\usepackage{titlesec}
\usepackage{enumitem}
\usepackage{booktabs}
\usepackage{fancyhdr}
\usepackage{graphicx}
\graphicspath{{figures/}{../../papers/arxiv-ac/figures/}}
\usepackage{ragged2e}
\usepackage{microtype}
\usepackage{eso-pic}
\usepackage{tikz}
\usepackage{hyperref}
\usepackage{xurl}

\definecolor{acpink}{RGB}{180,72,135}
\definecolor{acpurple}{RGB}{120,80,180}
\definecolor{acdark}{RGB}{64,56,74}
\definecolor{acgray}{RGB}{119,119,119}
\definecolor{palspink}{RGB}{205,92,155} % #cd5c9b — canonical purple-pals.svg fill

% === PALS SIDE-STAMPS (mirrors bios.mjs side-stamp pattern + ac-paper-layout.sty) ===
\AddToShipoutPictureBG{%
  \begin{tikzpicture}[remember picture, overlay]
    \pgfmathsetmacro{\palswiggle}{7*sin(\thepage * 67)}%
    \node[opacity=0.10, anchor=center, rotate={90 + \palswiggle}]
      at ([xshift=-1.0cm, yshift=-0.33\paperheight]current page.north west)
      {
\includegraphics[width=13cm]{palsmark}};
    \node[opacity=0.10, anchor=center, rotate={-90 + \palswiggle}]
      at ([xshift=1.0cm, yshift=-0.33\paperheight]current page.north east)
      {
\includegraphics[width=13cm]{palsmark}};
  \end{tikzpicture}%
}

\hypersetup{colorlinks=true,linkcolor=acpurple,urlcolor=acpurple,citecolor=acpurple,
  pdftitle={The Pals Mark: A History of the Aesthetic.Computer Logo}}

\titleformat{\section}{\normalfont\bfseries\normalsize\uppercase}{\thesection.}{0.5em}{}
\titlespacing{\section}{0pt}{1.2em}{0.3em}
\titleformat{\subsection}{\normalfont\bfseries\small}{\thesubsection}{0.5em}{}
\titlespacing{\subsection}{0pt}{0.8em}{0.2em}

\pagestyle{fancy}\fancyhf{}
\renewcommand{\headrulewidth}{0pt}
\fancyhead[C]{\footnotesize\color{acpink}\textit{The Pals Mark --- \ac{}}}
\fancyfoot[C]{\footnotesize\thepage}

\setlist[itemize]{nosep, leftmargin=1.2em, itemsep=0.15em}
\setlist[enumerate]{nosep, leftmargin=1.4em}
\setlength{\columnsep}{1.8em}
\setlength{\parindent}{1em}
\setlength{\parskip}{0.3em}

\newcommand{\ac}{\textsc{Aesthetic}{\color{acpink}.}\textsc{Computer}}
\newcommand{\acdot}{{\color{acpink}.}}
\newcommand{\pals}{\textsc{Pals}}
% commit link: \commit{shorthash}{message} -> links the hash to GitHub
\newcommand{\commit}[2]{\href{https://github.com/whistlegraph/aesthetic-computer/commit/#1}{\texttt{#1}}, \emph{``#2''}}
% repo file link
\newcommand{\repofile}[1]{\href{https://github.com/whistlegraph/aesthetic-computer/blob/main/#1}{\texttt{#1}}}

% Cover language bar. The same paper is published in six languages; each
% edition links to the others by canonical permalink. \klanghere marks the
% current language; \klang links a sibling edition by its URL suffix.
\newcommand{\klang}[2]{\href{https://papers.aesthetic.computer/pals-mark-26-arxiv#2.pdf}{\color{acgray}#1}}
\newcommand{\klanghere}[1]{{\bfseries\color{acpink}#1}}
\newcommand{\klangsep}{{\color{acgray!60}\,\textperiodcentered\,}}

\begin{document}

\twocolumn[{%
\begin{center}
{\acbold\fontsize{27pt}{31pt}\selectfont\color{acdark} The \pals{} Mark}\par
\vspace{0.5em}
\includegraphics[width=0.6\textwidth]{figures/cover}\par
\vspace{0.4em}
{\itshape\fontsize{14.5pt}{17pt}\selectfont\color{acpink} A History of the Aesthetic{\acdot}Computer Logo}\par
\vspace{0.5em}
{\small\color{acgray} \href{https://prompt.ac/@jeffrey}{@jeffrey} \,$\cdot$\, Aesthetic{\acdot}Computer \,$\cdot$\, logo by Niki Stebbins \,$\cdot$\, ORCID \href{https://orcid.org/0009-0007-4460-4913}{0009-0007-4460-4913}}\par
\vspace{0.45em}
{\footnotesize\sffamily\klanghere{EN}\klangsep\klang{DA}{-da}\klangsep\klang{ES}{-es}\klangsep\klang{ZH}{-zh}\klangsep\klang{JA}{-ja}\klangsep\klang{RU}{-ru}}\par
\vspace{0.55em}
{\color{acpink}\rule{\textwidth}{1.2pt}}
\vspace{0.5em}
\end{center}

\begin{quote}
\small\noindent\textbf{Abstract.}
The \pals{} are the logo of \ac{}: two hand-drawn figures holding hands with
their arms raised. This paper tells how the mark was made --- commissioned
from the artist Niki Stebbins and drawn, over the summer of 2023,
\emph{inside} \ac{} using its own \texttt{wgr} whistlegraph tool --- and how
it then became infrastructure: a rotating \texttt{pals.aesthetic.computer}
``logo grabber'' that serves the mark's many poses across avatars, papers,
native OS, and the web, each one a link back to the project.
\end{quote}
\vspace{0.5em}
}]

\section{What \ac{} is}

\ac{} is a mobile-first runtime and social network for creative computing
--- a single URL bar that doubles as a musical-instrument-like interface,
where memorizable commands summon interactive programs called
``pieces.'' It supports both JavaScript (\texttt{.mjs}) and a homegrown
Lisp dialect, \textbf{KidLisp} (118 built-ins). The whole system speaks to
pieces through one core API (\texttt{disk.mjs}, $\sim$600\,KB).

\textbf{When it went online.} The \texttt{aesthetic.computer} domain has
been held since \textbf{2013-12-12}. The platform as it exists was
conceived out of \emph{No Paint} (2020), development began in
\textbf{2021}, and it has deployed continuously from the current
repository since \textbf{December 2022} ($\sim$20{,}300 commits).

\textbf{Platform statistics} (live, \today):
\begin{itemize}
  \item $\sim$400 built-in pieces in-repo (383 \texttt{.mjs} + 20
        \texttt{.lisp}); 265 user-published works.
  \item \textbf{17{,}538} KidLisp programs stored in production --- a
        growing corpus of community-authored generative art.
  \item ATProto / Bluesky federation; per-handle record collections
        (paintings, moods, pieces, KidLisp) on the \texttt{at://} network.
\end{itemize}

\textbf{Who it serves, and who knows the mark.} \ac{}'s users are a
community of creative coders, artists, and non-technical ``Pals'' who
publish instantly under \texttt{@handle} pieces. The platform engages a
broad readership across the art-and-technology field, and \ac{} works have
entered institutional collections. The
\pals{} mark itself rides every one of these surfaces: the boot screen,
the papers, the social avatars, the native OS. Its conceptual lineage in
the radical-computer-music community of \emph{Goodiepal \& Pals} (see the
lineage section below) gives the name a recognized history in media art
beyond \ac{}'s own audience.

\section{What the mark depicts}

The \pals{} logo is \textbf{two outlined human figures holding hands with
their arms raised} --- like two friends walking or dancing together ---
rendered in a loose, organic, hand-drawn line. The canonical static
vector (\texttt{system/public/purple-pals.svg}) is pink/magenta,
\texttt{\#cd5c9b}. It is \emph{not} ``a pair of eyes,'' an earlier
misreading that has since been corrected in the paper stack.

Crucially, the mark was made \emph{with} the platform, not merely
\emph{for} it: Niki drew the logo explorations using \texttt{wgr}, the
whistlegraph-renderer drawing piece that runs inside \ac{} itself. The
logo is therefore both an identity and a demonstration of the tool.

\section{How the mark was made}

\subsection{A stamp, first}
It starts in the summer of \textbf{2020} with a small gift. Niki Stebbins
emails a drawing --- a little purple star --- and Jeffrey drops it into the
game as a stamp under the handle \texttt{@niki} and asks for more. For the
next two years the collaboration is a loose, ongoing \emph{drawing session}:
stamps, a drawing game, whistlegraph doodles. By \textbf{2022} it is paid
design work. None of it is the logo yet, but it is the soil the logo grows
from --- a habit of two people drawing together inside the tool.

\subsection{Reopening the design (June 2023)}
The logo proper begins with dissatisfaction. In \textbf{June 2023} Niki
writes that she is ``not 100\% satisfied with the logo I left you with,''
and reopens it. She thinks out loud about color --- ``like purple or
green\dots{} not blue, too many blue app logos'' --- and about form: should
the mark be ``like a frame, a closed shape only,'' something you could work
\emph{inside}? She sends a sketchbook page of candidates and picks a
favorite --- ``I like the circled one the most'' --- then immediately
worries it is ``pretty similar to are.na's logo,'' and keeps going. The key
instruction she gives herself is a correction of proportion: push it ``so it
looks less like stars and more like figures.'' A page of the studies gets a
verdict of its own --- ``looks like starfish lol'' --- and to fix it she
leaves the screen entirely, tying and photographing lengths of \textbf{rope}
in a park to understand how two linked bodies actually fall.

\subsection{Building the two figures (July--August 2023)}
Through the summer the mark resolves into \emph{two} figures rather than one.
Niki buys a used \textbf{Kodak EasyShare} digital camera specifically to
photograph the logo, sizes a version to ``the exact size of the VS
Marketplace logo,'' and asks whether Jeffrey wants ``more logo drawings in
a.c for social media pfps and the app-store thing.'' The repository agrees
with the messages: on \commit{5bcefc5a9}{better defaults in \texttt{wgr} for
niki's logo efforts} the drawing tool itself is tuned for her, and by
\commit{a0009f5d0}{add new vscode extension logo from niki} a shipped asset
lands. Late August she reports the ``variations of the \textbf{two logos}''
finished, describing the gesture precisely --- ``the right figure was
pulling at the left one\dots{} the left figure's arm is higher.'' The mark is
now the thing it still is: two people, arms up, one pulling the other along.
On \textbf{2023-08-31} an Instagram profile picture is chosen.

\subsection{Ice cream, and going live (September 2023)}
The finish is the strangest and best part. To photograph the mark in the
world, Niki draws it at full size on pavement and tries to shoot it with
ice cream --- but the ballpark ice cream comes in cups, so she detours to a
McDonald's for a cone and gets the shot. On \textbf{2023-09-09} she ``sent
the logos to a.c email'' under the subject \emph{!c3 cr34m l0go\$} (body:
``ice cream logos 8)''). That same week the mark stops being a picture and becomes
infrastructure: \commit{0aba1a378}{add dynamic logo generator} lands, and
Jeffrey reports back the one-line result --- ``\^{}\ will now return a random
logo every time.'' The logo was drawn \emph{with} the platform (in the
\texttt{wgr} whistlegraph tool) and then served \emph{by} it. It was never
only a picture; it was always also a program.

\section{The endpoint: \texttt{pals.aesthetic.computer}}

The mark became infrastructure in September 2023, when \ac{} stood up a
dedicated logo service and gave it a memorizable home.

\begin{itemize}
  \item \textbf{2023-09-05} --- \commit{0aba1a378}{add dynamic logo
        generator}. The backend module \repofile{system/backend/logo.mjs}
        is stamped \texttt{23.09.05.22.31}.
  \item \textbf{2023-09-12} --- \commit{f1a4f272c}{logo -> pals.aesthetic}
        --- the route is renamed to the \pals{} identity.
  \item \textbf{2023-09-13} --- \commit{12e348035}{pals.aesthetic.computer
        url fix} --- the endpoint is live and addressable at its own
        subdomain.
\end{itemize}

\noindent So \href{https://pals.aesthetic.computer}{\texttt{pals.aesthetic.computer}}
\textbf{went online on 2023-09-12--13}, three days after the dynamic
generator landed.

\subsection{How it works as a ``logo grabber''}

The endpoint is a tiny rotating proxy, not a single fixed image. Each
request returns a different \pals{} pose:

\begin{enumerate}
  \item \href{https://pals.aesthetic.computer}{\texttt{pals.aesthetic.computer}}
        redirects (via \repofile{system/netlify.toml}) to
        \href{https://aesthetic.computer/api/logo}{\texttt{aesthetic.computer/api/logo}}.
  \item \repofile{system/backend/logo.mjs} holds a list of \pals{}
        \emph{paintings} on the \texttt{pals-aesthetic-computer} CDN
        bucket and returns a random one each call (avoiding an immediate
        repeat).
  \item The function content-negotiates: a browser gets a full-bleed HTML
        page on a random purple, tappable to cycle to the next pose
        (passing \texttt{?previousLogo=} so it never repeats); a
        \texttt{curl}/\texttt{<img>}/server request gets the raw
        \texttt{.png} directly. The same URL is thus both a viewer and a
        hot-linkable asset.
\end{enumerate}

\subsection{The many poses of the \pals{}}

The grabber's pose set is a gallery of \texttt{wgr} paintings from
mid-2023 (the \texttt{painting-2023.7.29}--\texttt{2023.8.21} series).
Because the logo is a \emph{collection of drawn poses} rather than one
frozen glyph, the brand reads as alive: every page load, avatar refresh,
or favicon fetch can show the two figures in a slightly different stance.
The full set of twelve live poses is reproduced in Appendix~A.

\section{Use across properties}

Because any property can hot-link
\href{https://pals.aesthetic.computer}{\texttt{pals.aesthetic.computer}}
(or embed the static
\repofile{system/public/purple-pals.svg}), a single mark threads the
whole ecosystem, and every instance is a link back to the project:

\begin{itemize}
  \item \textbf{Social avatars \& Gravatar} --- the rotating endpoint
        supplies a profile image that quietly changes pose, so the same
        identity appears across services while always resolving to
        \ac{}.
  \item \textbf{Web \& marketing} --- favicons, docs headers, and
        embedded \texttt{<img>} tags pull the live \pals{}; a tap cycles
        poses, turning the logo itself into a small interactive teaser
        for the platform.
  \item \textbf{Papers} --- the \pals{} ride the margins of every \ac{}
        working draft as semi-opaque, page-wiggling side-stamps (see this
        page), and appear in print-edition colophons.
  \item \textbf{Native OS \& desktop} --- the \pals{} serve as boot
        branding (Plymouth splash) in AC Native OS and as the vector
        tray icon in AC Electron and MenuBand --- accent-tinted, badge-able.
  \item \textbf{Print (pending)} --- in May 2025 Jeffrey contacted Casey
        Rubber Stamps to cut physical \pals{} stamps in three sizes;
        the order stalled and was never completed.
\end{itemize}

\section{Why ``\pals{}''? The Goodiepal lineage}

The name is not decoration; it is a structural claim, and its weight
comes from outside \ac{}. As argued in the \ac{} paper
\emph{Radical Computer Art: Goodiepalian Approaches in Aesthetic
Computer} (\repofile{papers/arxiv-goodiepal/goodiepal.tex}) and the
platform paper it cites, ``\pals{}'' descends from \textbf{Goodiepal \&
Pals} --- the loose affinity group around the Danish-Faroese musician and
Radical Computer Music theorist Goodiepal (Gaeoudji SYGNOK, b.~1974).

The \pals{} model is a politics of organization: a group with no
membership, no institutional backing, and no formal structure, held
together by shared practice and mutual recognition rather than
affiliation. Goodiepal taught History and Aesthetics of Electronic Music
at the Royal Danish Academy of Music and left in 2008 after a conflict
described as having ``the wrong type of friends,'' having declared
``intellectual war against the stupidity in modern computer music and
media art.'' Jeffrey toured with Goodiepal \& Pals across Europe in 2018,
and the tools built for those performances are direct ancestors of
\ac{}'s architecture.

So the logo carries that media-arts political history into the brand. As
the Goodiepal paper puts it, naming the logo ``\pals{}'' --- and titling
the score \emph{Score for Aesthetic{\color{acpink}.}Computer \& Pals} --- ``is not
branding; it is a structural claim: the platform belongs to its community
of practitioners, not to its developer or its investors.'' Two figures
holding hands, made \emph{inside} the tool, is the argument rendered as a
mark: \ac{} is \emph{for} anyone who shows up --- for \pals{}.

\section{Authorship}

Jeffrey hired \textbf{Niki Stebbins} to design the \ac{} logo under his
direction. Across 2022--2023 she developed the mark inside the platform's
own \texttt{wgr} drawing tool, working through the sketchbook, rope, and
ice-cream explorations documented above and delivering the two-figures
device that ships as \texttt{purple-pals.svg}. The rotating
\texttt{pals.aesthetic.computer} grabber went live on 2023-09-12--13; the
canonical static vector first appears on 2024-01-31
(\commit{bf29eadf8}{host purple-pals graphic for docs}); and from 2026-02
the repository and papers explicitly center ``\pals{}'' as platform
identity (\commit{48dd8923e}{rename score to ``Score for
Aesthetic.Computer \& Pals''}). The mark is a commissioned work: authored
by Niki, directed by Jeffrey, and owned by \ac{}.

\vspace{0.6em}
{\footnotesize\color{acgray}\noindent Sources: repository commit history;
\repofile{system/backend/logo.mjs};
\repofile{system/netlify/functions/logo.js};
\repofile{system/netlify.toml};
\repofile{papers/arxiv-goodiepal/goodiepal.tex};
\repofile{system/public/purple-pals.svg}.\par}

% ============================ APPENDICES ============================
\onecolumn
\appendix

\noindent\emph{Note on dates.} Every visual asset below is captioned with
the date it was made: \texttt{wgr} poses by the timestamp encoded in the
painting slug; monorepo assets by their first-commit (git
\texttt{--diff-filter=A}) date; live captures by the date they were
photographed from production.

\section*{Appendix A --- The many poses (live \texttt{pals.aesthetic.computer} set)}
\addcontentsline{toc}{section}{Appendix A --- Poses}
\noindent The twelve \texttt{wgr} paintings the logo grabber rotates
through, pulled live from the \texttt{pals-aesthetic-computer} CDN. Same
two-figure motif, each a different stance and color treatment.
\vspace{1.2em}

\begin{minipage}[t]{0.242\textwidth}\centering
  
\includegraphics[width=\linewidth,height=\linewidth,keepaspectratio]{figures/poses/small/pose-1.png}\\[2pt]
  {\scriptsize\textbf{pose 1}}\\[-1pt]
  {\scriptsize\color{acgray}2023-08-01}
\end{minipage}\hfill\begin{minipage}[t]{0.242\textwidth}\centering
  \includegraphics[width=\linewidth,height=\linewidth,keepaspectratio]{figures/poses/small/pose-2.png}\\[2pt]
  {\scriptsize\textbf{pose 2}}\\[-1pt]
  {\scriptsize\color{acgray}2023-08-21}
\end{minipage}\hfill\begin{minipage}[t]{0.242\textwidth}\centering
  
\includegraphics[width=\linewidth,height=\linewidth,keepaspectratio]{figures/poses/small/pose-3.png}\\[2pt]
  {\scriptsize\textbf{pose 3}}\\[-1pt]
  {\scriptsize\color{acgray}2023-08-21}
\end{minipage}\hfill\begin{minipage}[t]{0.242\textwidth}\centering
  \includegraphics[width=\linewidth,height=\linewidth,keepaspectratio]{figures/poses/small/pose-4.png}\\[2pt]
  {\scriptsize\textbf{pose 4}}\\[-1pt]
  {\scriptsize\color{acgray}2023-08-01}
\end{minipage}

\vspace{0.8em}

\begin{minipage}[t]{0.242\textwidth}\centering
  \includegraphics[width=\linewidth,height=\linewidth,keepaspectratio]{figures/poses/small/pose-5.png}\\[2pt]
  {\scriptsize\textbf{pose 5}}\\[-1pt]
  {\scriptsize\color{acgray}2023-08-21}
\end{minipage}\hfill\begin{minipage}[t]{0.242\textwidth}\centering
  \includegraphics[width=\linewidth,height=\linewidth,keepaspectratio]{figures/poses/small/pose-6.png}\\[2pt]
  {\scriptsize\textbf{pose 6}}\\[-1pt]
  {\scriptsize\color{acgray}2023-08-01}
\end{minipage}\hfill\begin{minipage}[t]{0.242\textwidth}\centering
  \includegraphics[width=\linewidth,height=\linewidth,keepaspectratio]{figures/poses/small/pose-7.png}\\[2pt]
  {\scriptsize\textbf{pose 7}}\\[-1pt]
  {\scriptsize\color{acgray}2023-08-21}
\end{minipage}\hfill\begin{minipage}[t]{0.242\textwidth}\centering
  \includegraphics[width=\linewidth,height=\linewidth,keepaspectratio]{figures/poses/small/pose-8.png}\\[2pt]
  {\scriptsize\textbf{pose 8}}\\[-1pt]
  {\scriptsize\color{acgray}2023-08-21}
\end{minipage}

\vspace{0.8em}

\begin{minipage}[t]{0.242\textwidth}\centering
  \includegraphics[width=\linewidth,height=\linewidth,keepaspectratio]{figures/poses/small/pose-9.png}\\[2pt]
  {\scriptsize\textbf{pose 9}}\\[-1pt]
  {\scriptsize\color{acgray}2023-08-01}
\end{minipage}\hfill\begin{minipage}[t]{0.242\textwidth}\centering
  \includegraphics[width=\linewidth,height=\linewidth,keepaspectratio]{figures/poses/small/pose-10.png}\\[2pt]
  {\scriptsize\textbf{pose 10}}\\[-1pt]
  {\scriptsize\color{acgray}2023-08-01}
\end{minipage}\hfill\begin{minipage}[t]{0.242\textwidth}\centering
  \includegraphics[width=\linewidth,height=\linewidth,keepaspectratio]{figures/poses/small/pose-11.png}\\[2pt]
  {\scriptsize\textbf{pose 11}}\\[-1pt]
  {\scriptsize\color{acgray}2023-07-31}
\end{minipage}\hfill\begin{minipage}[t]{0.242\textwidth}\centering
  
\includegraphics[width=\linewidth,height=\linewidth,keepaspectratio]{figures/poses/small/pose-12.png}\\[2pt]
  {\scriptsize\textbf{pose 12}}\\[-1pt]
  {\scriptsize\color{acgray}2023-07-29}
\end{minipage}

\vspace{0.8em}




\newpage
\section*{Appendix B --- The mark across the monorepo}
\addcontentsline{toc}{section}{Appendix B --- Monorepo assets}
\noindent The same \pals{} drawn once and re-skinned everywhere: the
canonical vector and its descendants across the codebase --- editor
extension, native-OS boot splash, desktop tray, pop covers, invoice stamp,
and the keymaps colorway set. Each tile is captioned with its first-commit date.
\vspace{1.2em}

\begin{minipage}[t]{0.194\textwidth}\centering
  
\includegraphics[width=\linewidth,height=\linewidth,keepaspectratio]{figures/assets/c01-purple-pals.png}\\[2pt]
  {\scriptsize\textbf{purple-pals.svg}}\\[-1pt]
  {\scriptsize\color{acgray}2024-01-31}
\end{minipage}\hfill\begin{minipage}[t]{0.194\textwidth}\centering
  
\includegraphics[width=\linewidth,height=\linewidth,keepaspectratio]{figures/assets/c02-vscode.png}\\[2pt]
  {\scriptsize\textbf{vscode-extension}}\\[-1pt]
  {\scriptsize\color{acgray}2024-01-31}
\end{minipage}\hfill\begin{minipage}[t]{0.194\textwidth}\centering
  
\includegraphics[width=\linewidth,height=\linewidth,keepaspectratio]{figures/assets/c03-tvos.png}\\[2pt]
  {\scriptsize\textbf{apple/tvos-tapes}}\\[-1pt]
  {\scriptsize\color{acgray}2024-01-31}
\end{minipage}\hfill\begin{minipage}[t]{0.194\textwidth}\centering
  
\includegraphics[width=\linewidth,height=\linewidth,keepaspectratio]{figures/assets/c04-plymouth.png}\\[2pt]
  {\scriptsize\textbf{AC-OS boot splash}}\\[-1pt]
  {\scriptsize\color{acgray}2026-02-21}
\end{minipage}\hfill\begin{minipage}[t]{0.194\textwidth}\centering
  
\includegraphics[width=\linewidth,height=\linewidth,keepaspectratio]{figures/assets/c05-tray.png}\\[2pt]
  {\scriptsize\textbf{Electron tray}}\\[-1pt]
  {\scriptsize\color{acgray}2026-05-20}
\end{minipage}

\vspace{0.8em}

\begin{minipage}[t]{0.194\textwidth}\centering
  
\includegraphics[width=\linewidth,height=\linewidth,keepaspectratio]{figures/assets/c14-stamp.png}\\[2pt]
  {\scriptsize\textbf{invoice stamp}}\\[-1pt]
  {\scriptsize\color{acgray}2026-03-06}
\end{minipage}\hfill\begin{minipage}[t]{0.194\textwidth}\centering
  
\includegraphics[width=\linewidth,height=\linewidth,keepaspectratio]{figures/assets/c12-gray.png}\\[2pt]
  {\scriptsize\textbf{grayscale (quotes)}}\\[-1pt]
  {\scriptsize\color{acgray}2026-05-21}
\end{minipage}\hfill\begin{minipage}[t]{0.194\textwidth}\centering
  
\includegraphics[width=\linewidth,height=\linewidth,keepaspectratio]{figures/assets/c13-hellsine.png}\\[2pt]
  {\scriptsize\textbf{pop cover glyph}}\\[-1pt]
  {\scriptsize\color{acgray}2026-05-24}
\end{minipage}\hfill\begin{minipage}[t]{0.194\textwidth}\centering
  
\includegraphics[width=\linewidth,height=\linewidth,keepaspectratio]{figures/assets/c06-pink.png}\\[2pt]
  {\scriptsize\textbf{keymaps pink}}\\[-1pt]
  {\scriptsize\color{acgray}2026-06-13}
\end{minipage}\hfill\begin{minipage}[t]{0.194\textwidth}\centering
  
\includegraphics[width=\linewidth,height=\linewidth,keepaspectratio]{figures/assets/c07-blue.png}\\[2pt]
  {\scriptsize\textbf{keymaps blue}}\\[-1pt]
  {\scriptsize\color{acgray}2026-06-13}
\end{minipage}

\vspace{0.8em}

\begin{minipage}[t]{0.194\textwidth}\centering
  
\includegraphics[width=\linewidth,height=\linewidth,keepaspectratio]{figures/assets/c08-green.png}\\[2pt]
  {\scriptsize\textbf{keymaps green}}\\[-1pt]
  {\scriptsize\color{acgray}2026-06-13}
\end{minipage}\hfill\begin{minipage}[t]{0.194\textwidth}\centering
  
\includegraphics[width=\linewidth,height=\linewidth,keepaspectratio]{figures/assets/c09-orange.png}\\[2pt]
  {\scriptsize\textbf{keymaps orange}}\\[-1pt]
  {\scriptsize\color{acgray}2026-06-13}
\end{minipage}\hfill\begin{minipage}[t]{0.194\textwidth}\centering
  
\includegraphics[width=\linewidth,height=\linewidth,keepaspectratio]{figures/assets/c10-purple.png}\\[2pt]
  {\scriptsize\textbf{keymaps purple}}\\[-1pt]
  {\scriptsize\color{acgray}2026-06-13}
\end{minipage}\hfill\begin{minipage}[t]{0.194\textwidth}\centering
  
\includegraphics[width=\linewidth,height=\linewidth,keepaspectratio]{figures/assets/c11-yellow.png}\\[2pt]
  {\scriptsize\textbf{keymaps yellow}}\\[-1pt]
  {\scriptsize\color{acgray}2026-06-13}
\end{minipage}\hfill

\newpage
\section*{Appendix C --- Live placements (real UIs)}
\addcontentsline{toc}{section}{Appendix C --- Live placements}
\noindent The \pals{} photographed living in real production interfaces,
each captioned with its source URL (links are live): the Visual Studio
Marketplace listing (the \pals{} as the extension's product icon), the
monorepo README on both GitHub and the canonical Tangled host, and the
rotating \href{https://pals.aesthetic.computer}{\texttt{pals.aesthetic.computer}}
endpoint (same URL, different poses on reload --- plus the raw
\texttt{/api/logo.png} handed to \texttt{<img>} tags and avatars).
\vspace{1.0em}

% Hand-authored: the Pals living in real production UIs, each with a caption + live URL.

\newcommand{\uishot}[3]{%
  % #1 image  #2 caption  #3 url
  \begin{center}
  \fbox{\includegraphics[width=0.96\linewidth]{#1}}\\[0.4em]
  {\footnotesize #2}\\[0.1em]
  {\footnotesize\ttfamily\href{#3}{#3}}
  \end{center}
  \vspace{0.7em}
}
\newcommand{\uishottwo}[6]{%
  % (#1 img #2 cap #3 url) (#4 img #5 cap #6 url)
  \noindent
  \begin{minipage}[t]{0.485\linewidth}\centering
    \fbox{\includegraphics[width=0.97\linewidth]{#1}}\\[0.35em]
    {\footnotesize #2}\\[0.1em]{\scriptsize\ttfamily\href{#3}{#3}}
  \end{minipage}\hfill
  \begin{minipage}[t]{0.485\linewidth}\centering
    \fbox{\includegraphics[width=0.97\linewidth]{#4}}\\[0.35em]
    {\footnotesize #5}\\[0.1em]{\scriptsize\ttfamily\href{#6}{#6}}
  \end{minipage}
  \vspace{0.9em}
}

\uishot{figures/live/ui-vscode-marketplace.png}{%
  \textbf{Visual Studio Marketplace} --- the \pals{} as the product icon of
  the official \ac{} VS~Code extension (1{,}201 installs).}{%
  https://marketplace.visualstudio.com/items?itemName=aesthetic-computer.aesthetic-computer-code}

\uishottwo
  {figures/live/ui-github-readme.png}{\textbf{GitHub} --- the \pals{} atop the
   monorepo README.}{https://github.com/whistlegraph/aesthetic-computer}
  {figures/live/ui-tangled-readme.png}{\textbf{Tangled} (canonical git host)
   --- the same README, same mark.}{https://tangled.org/aesthetic.computer/core}

\uishot{figures/compare/pals-live-strip.png}{%
  \textbf{\texttt{pals.aesthetic.computer}} in a live browser --- the same URL
  reloaded shows different poses (the rotating ``logo grabber''); the raw
  \texttt{/api/logo.png} is what \texttt{<img>} tags and avatars (Gravatar,
  social) fetch across the web.}{%
  https://pals.aesthetic.computer}



\end{document}
